% ============================================================
%  Worked example: estimability in a two-receiver network
%
%  A fully numerical, noise-free demonstration of the rank
%  deficiency / minimum-constraint datum / lumping material in
%  the Network-mode section (\dcref{sec:net-sbasis}{the minimum-constraint datum section of the Network-mode document}).  Smallest
%  non-trivial case: two receivers, four satellites, code-only
%  clocks first, then carrier phase to show why the satellite
%  phase bias is estimable only across stations.
%
%  Standalone-compilable subfile (shares gnssreport.sty via the
%  subfiles package, as network.tex does).  Wired into the
%  combined volume after \subfile{network} in main.tex.
% ============================================================
\documentclass[main.tex]{subfiles}
\begin{document}
\ifSubfilesClassLoaded{\begin{center}\small\itshape Part of the GnssProcessor Mathematical Reference \textbullet\ \today\end{center}}{}

\section{Worked example: estimability in a two-receiver network}
\label{sec:net-example}

This section works the rank deficiency and minimum-constraint datum of
\dcref{sec:net-sbasis}{the minimum-constraint datum section of the Network-mode document}
through numerically on the smallest network that shows the structure: two
receivers $\rxset = \{A, B\}$ and four satellites $\sat \in \{1,2,3,4\}$.
Everything not needed for the estimability argument is removed. Positions
are known, so the geometric range drops out; a single frequency is used;
the troposphere and ionosphere are absent; the data are noise-free. Every
number below is therefore exact, and the gauge freedoms are visible by
inspection. The subject of the example is the null space of the system,
not its size.

\subsection{Code only: clocks and one gauge freedom}
\label{sec:net-example-code}

With geometry removed, the single-frequency pseudorange of
\dcref{sec:obs}{the observation-equations section} collapses to a
difference of two clock lumps,
\begin{equation}
\code^\sat_\rx \;=\; a_\rx \;-\; g^\sat,
\qquad
a_\rx := \sLight\,\mathrm{d}t_\rx + B_\rx,
\quad
g^\sat := \sLight\,\mathrm{d}t^\sat + B^\sat,
\label{eq:ex-code-model}
\end{equation}
where $a_\rx$ is the receiver clock and $g^\sat$ the satellite clock. On
a single frequency the receiver and satellite \emph{code biases}
$B_\rx, B^\sat$ are inseparable from the respective clocks, so each
``clock'' here is a clock\,+\,code-bias lump. With one band, the
$\mathrm{IF}$-vs-$\mathrm{DCB}$ split of
\dcref{sec:net-phasebias}{the satellite phase-bias section of the Network-mode document}
does not arise. The system has $|\rxset| + |\{\sat\}| = 2 + 4 = 6$
unknowns and $2 \times 4 = 8$ observations.

The ground truth (metres) is
\begin{equation*}
a_A = 5,\quad a_B = 3,\qquad
g^1 = 2,\quad g^2 = -1,\quad g^3 = 4,\quad g^4 = 0.5,
\end{equation*}
and \Cref{eq:ex-code-model} maps it to the eight pseudoranges
\begin{center}
\begin{tabular}{@{}lcccc@{}}
\toprule
$\code^\sat_\rx$ & $\sat=1$ & $\sat=2$ & $\sat=3$ & $\sat=4$ \\
\midrule
$\rx = A$ & $3$ & $6$ & $1$ & $4.5$ \\
$\rx = B$ & $1$ & $4$ & $-1$ & $2.5$ \\
\bottomrule
\end{tabular}
\end{center}

\paragraph{The deficiency from two directions.}
\emph{(i) Invariance.} Adding a constant $\kappa$ to all six unknowns,
$a_\rx \mapsto a_\rx + \kappa$ and $g^\sat \mapsto g^\sat + \kappa$,
leaves every difference $a_\rx - g^\sat$, and therefore every
observation, unchanged. The truth above and the shifted set
$\{15,13,12,9,14,10.5\}$ fit the data identically. The data determine
clock differences only, never absolute clocks.

\emph{(ii) Singular normals.} With the unknowns ordered as
$\bm{x} = [\,a_A, a_B, g^1, g^2, g^3, g^4\,]^\Tr$ and unit weights,
the normal matrix $\bm{N} = \bm{A}^\Tr\bm{A}$ is
\begin{equation}
\bm{N} =
\begin{bmatrix}
\phantom{-}4 & \phantom{-}0 & -1 & -1 & -1 & -1 \\
\phantom{-}0 & \phantom{-}4 & -1 & -1 & -1 & -1 \\
-1 & -1 & \phantom{-}2 & \phantom{-}0 & \phantom{-}0 & \phantom{-}0 \\
-1 & -1 & \phantom{-}0 & \phantom{-}2 & \phantom{-}0 & \phantom{-}0 \\
-1 & -1 & \phantom{-}0 & \phantom{-}0 & \phantom{-}2 & \phantom{-}0 \\
-1 & -1 & \phantom{-}0 & \phantom{-}0 & \phantom{-}0 & \phantom{-}2
\end{bmatrix},
\qquad
\bm{N}\,\bm{1} = \bm{0},
\label{eq:ex-code-normals}
\end{equation}
a $6\times6$ matrix of rank $5$. The all-ones vector
$\bm{1} = [1,1,1,1,1,1]^\Tr$ is an exact null vector (every row sums to
zero), and it is the common-shift direction of view~(i). The system
$\bm{N}\bm{x} = \bm{u}$ has no unique solution without one added
constraint, a \emph{datum}.

\paragraph{Two datum choices, same fit.}
Pinning one clock removes the gauge. Two natural choices give different
parameter values, and both fit the eight observations exactly:
\begin{center}
\begin{tabular}{@{}lcccccc@{}}
\toprule
datum (S-basis) & $a_A$ & $a_B$ & $g^1$ & $g^2$ & $g^3$ & $g^4$ \\
\midrule
$a_A \equiv 0$ (reference receiver) & $0$ & $-2$ & $-3$ & $-6$ & $-1$ & $-4.5$ \\
$g^1 \equiv 0$ (reference satellite) & $3$ & $\phantom{-}1$ & $\phantom{-}0$ & $-3$ & $\phantom{-}2$ & $-1.5$ \\
\bottomrule
\end{tabular}
\end{center}
Under $a_A \equiv 0$, every satellite-clock estimate comes out as
$\hat g^\sat = g^\sat - a_A^{\text{true}}$: the reference receiver's
clock is absorbed into all of them. The estimated ``satellite clock'' is
the physical clock shifted by the chosen datum. This offset is the
\emph{lumping}, and the user must account for it. For this reason a
producer declares its datum
(\dcref{sec:net-closedloop}{the closed-loop section of the Network-mode document}).

\paragraph{What is invariant.}
The estimable quantities (combinations annihilated by the null vector
$\bm{1}$, i.e.\ differences) are identical in both bases:
\begin{center}
\begin{tabular}{@{}lccc@{}}
\toprule
estimable quantity & $a_A \equiv 0$ & $g^1 \equiv 0$ & (truth) \\
\midrule
$a_A - a_B$  & $2$  & $2$  & $2$ \\
$g^2 - g^1$  & $-3$ & $-3$ & $-3$ \\
$g^3 - g^1$  & $2$  & $2$  & $2$ \\
$a_B - g^4$  & $2.5$ & $2.5$ & ($=\code^4_B$) \\
\bottomrule
\end{tabular}
\end{center}
The absolute values are a convention; the differences are physical.
This is the S-system of \cite{teunissen1985generalised} in six
parameters.

\begin{remark}[Why the producer cannot difference]
The double difference of this code-only model is
$\dd{\code^{\sat,1}_{AB}} = (\code^1_A - \code^\sat_A) - (\code^1_B -
\code^\sat_B)$. The between-satellite step gives $g^\sat - g^1$ for both
receivers, so the between-receiver step yields exactly $\bm{0}$. Double
differencing eliminates the entire model: DD removes clocks, and here the
model contains nothing else
(\dcref{sec:diff}{the differencing section}). A DD user estimates
nothing. The network producer instead keeps the clocks and pins one
datum. This degeneracy is, in miniature, the reason a network producer
exists.
\end{remark}

\subsection{Adding carrier phase: the satellite phase bias needs the network}
\label{sec:net-example-phase}

The single-frequency carrier phase, with the same clock lumps, is
\begin{equation}
\phaseo^\sat_\rx \;=\; a_\rx - g^\sat \;+\; b^\sat \;+\; N^\sat_\rx,
\label{eq:ex-phase-model}
\end{equation}
where $b^\sat$ is the satellite phase bias, the producer's target,
and $N^\sat_\rx$ is the float ambiguity (in metres). The receiver phase
bias and the receiver code-minus-phase offset are constant per receiver
and are absorbed into $N^\sat_\rx$. The satellite bias $b^\sat$ stays
explicit because the producer emits it. The truth is extended with
\begin{align*}
b^1 &= 0.10, & b^2 &= -0.20, & b^3 &= 0.05, & b^4 &= 0.15, \\
N^\sat_A &: 10,\ -5,\ 20,\ 8, & & & N^\sat_B &: 12,\ -3,\ 18,\ 9, & &
\end{align*}
giving the eight carrier-phase observations
\begin{center}
\begin{tabular}{@{}lcccc@{}}
\toprule
$\phaseo^\sat_\rx$ & $\sat=1$ & $\sat=2$ & $\sat=3$ & $\sat=4$ \\
\midrule
$\rx = A$ & $13.10$ & $0.80$ & $21.05$ & $12.65$ \\
$\rx = B$ & $13.10$ & $0.80$ & $17.05$ & $11.65$ \\
\bottomrule
\end{tabular}
\end{center}

\paragraph{Counting: one receiver vs.\ two.}
The code fixes the clocks, so per satellite the phase contributes
$b^\sat + N^\sat_\rx$. With one receiver this is a single equation
$b^\sat + N^\sat_A$ in two unknowns. The satellite phase bias is then
inseparable from that receiver's ambiguity, and the equation carries no
redundancy. With two receivers there are two equations,
\begin{equation}
b^\sat + N^\sat_A, \qquad b^\sat + N^\sat_B,
\label{eq:ex-phase-pair}
\end{equation}
in three unknowns $(b^\sat, N^\sat_A, N^\sat_B)$: rank $2$, one
deficiency per satellite. The estimable two-dimensional space contains
the \emph{between-receiver single-difference ambiguity}
$\brd{N^\sat_{AB}} = N^\sat_A - N^\sat_B$. This quantity does not exist
with one receiver, and it is the integer a user resolves
(\dcref{sec:diff}{the differencing section}). The second receiver gives
$b^\sat$ an identity shared across stations.

\paragraph{Datum and solution.}
The joint code\,+\,phase system has a $1$-dimensional clock gauge (as
before) plus one phase-bias$\leftrightarrow$ambiguity gauge per
satellite: $b^\sat \mapsto b^\sat + c_\sat$, $N^\sat_\rx \mapsto
N^\sat_\rx - c_\sat$ for all $\rx$. The minimum-constraint datum used
here has two parts. The first is the clock pin $a_A \equiv 0$. It is
S-transform-equivalent to the reference-\emph{satellite} clock pin
$g^1 \equiv 0$ that the estimator's S-basis declares
(\dcref{sec:net-sbasis}{the minimum-constraint datum section of the Network-mode document}):
in its default pin mode the shared-block gauge pins one satellite clock
per constellation, not a receiver clock (the optional
\texttt{network.clock\_datum\_sum\_to\_zero} and
\texttt{network.clock\_datum\_broadcast\_anchored} modes replace the
pin). The second part is the set of reference-receiver ambiguity pins
$N^\sat_A \equiv 0$ for every $\sat$: the local datum on the pivot
receiver, folded into station $A$'s rows. The solution is
\begin{center}
\begin{tabular}{@{}lcccc@{}}
\toprule
& $\sat=1$ & $\sat=2$ & $\sat=3$ & $\sat=4$ \\
\midrule
emitted phase bias $\bar b^\sat = b^\sat + N^\sat_A$ & $10.10$ & $-5.20$ & $20.05$ & $8.15$ \\
$\hat N^\sat_A$ (datum) & $0$ & $0$ & $0$ & $0$ \\
$\hat N^\sat_B = N^\sat_B - N^\sat_A$ & $2$ & $2$ & $-2$ & $1$ \\
\bottomrule
\end{tabular}
\end{center}
with the clocks unchanged from \Cref{sec:net-example-code}. The fit
checks directly, e.g.\
$\hat{\phaseo}^3_B = a_B - g^3 + \bar b^3 + \hat N^3_B
   = -2 -(-1) + 20.05 + (-2) = 17.05$, matching the observation.

\paragraph{Reading the result.}
The solution shows two properties.
\begin{itemize}[leftmargin=*, itemsep=1pt]
\item \emph{The lump dominates the emitted bias.} The physical
      satellite phase bias of satellite~$1$ is $b^1 = 0.10$~m, while the
      emitted estimable bias is $\bar b^1 = 10.10$~m: it has absorbed
      receiver~$A$'s ambiguity ($10$~m) for that satellite. This is
      Odijk's $\bar b^\sat = b^\sat + (\text{pivot-receiver
      ambiguity})$ (\dcref{sec:net-phasebias}{the satellite phase-bias section of the Network-mode document},
      Eq.~4/9) in numbers. A consumer recovers integers only by removing
      the same lumped quantity, so the producer datum and the consumer
      integer recovery must belong to one S-system.
\item \emph{The network makes the integer estimable.} The estimable
      $\hat N^\sat_B = N^\sat_B - N^\sat_A = \brd{N^\sat_{AB}}$ is the
      between-receiver single-difference ambiguity, the quantity an RTK
      user fixes to an integer. It is absent from a single-receiver
      problem; the second station introduces it. The satellite phase
      bias becomes a definite, transferable number at the cost of
      pinning one reference ambiguity, and it carries that pin.
\end{itemize}

Up to the geometry and the second frequency it omits, this
four-satellite, two-receiver instance has the structure of the synthetic
network on which the estimator core is tested. The
\texttt{RecoversSyntheticTruthUpToSTransform} test
(\texttt{tests/test\_network\_estimator.cpp}) builds rows of this shape
(four receivers and five satellites, from a seeded random truth), pins
the reference-satellite clock datum $g^1 \equiv 0$
(S-transform-equivalent to the $a_A \equiv 0$ used above) through
\texttt{NetworkEstimator::solve}, and pins the pivot receiver's
ambiguities with tight-prior rows in that station's block. It asserts
recovery of the satellite-side parameters up to the declared
S-transform: the lumped $\bar b^\sat$ and the single-difference
ambiguities above, not the unobservable physical $b^\sat$.

\ifSubfilesClassLoaded{\printbibliography}{}
\end{document}
